% Generated draft: requires mathematical and source review.
\documentclass[11pt]{article}
\usepackage[margin=25mm]{geometry}
\usepackage{fontspec}
\setmainfont{Latin Modern Roman}
\usepackage{amsmath,amssymb,mathtools,unicode-math}
\setmathfont{Latin Modern Math}
\usepackage{enumitem,needspace,xurl,hyperref}
\setlength{\parindent}{0pt}
\setlength{\parskip}{5pt}
\setlength{\emergencystretch}{3em}
\newcommand{\N}{\mathbb{N}}
\newcommand{\Z}{\mathbb{Z}}
\newcommand{\Q}{\mathbb{Q}}
\newcommand{\R}{\mathbb{R}}
\newcommand{\C}{\mathbb{C}}
\newcommand{\F}{\mathbb{F}}
\newcommand{\A}{\mathbb{A}}
\newcommand{\PP}{\mathbb{P}}
\newcommand{\E}{\mathbb{E}}
\newcommand{\eps}{\varepsilon}
\newcommand{\dd}{\,\mathrm{d}}
\begin{document}
\section*{UnsolvedMath ID: 20000247}
\textbf{Source problem code:} catalogue entry
\subsection{Definitions and mathematical statement}
The reconstructed question, attributed to Anand Sawant, asks: when are the preceding Questions 1--9 on projective modules and \(\A^1\)-homotopy well-defined for singular varieties, where have they been considered, and what is their status? Those nine questions are not included in this record, so a complete source-specific status survey cannot be formulated from it. Different extensions of vector bundles, Chern classes, Chow groups and motivic classification can yield different singular versions.

The attached aid studies a free-summand part of that question. For a commutative noetherian ring \(A\) and finite projective \(A\)-module \(P\), write
\[
\operatorname{Um}(P^\vee)=\{\varphi\in\operatorname{Hom}_A(P,A):\varphi(P)=A\}.
\]
Such a functional splits and is equivalent to a decomposition \(P\simeq A\oplus Q\). The aid reports that for a nilpotent ideal \(J\subset A\), existence of such a free line summand is equivalent for \(P\) and \(P/JP\). This is existence invariance, not a bijection of all individual lifted functionals.

For reduced noetherian \(A\) with finite normalization \(B\), let \(I=\{a\in A:aB\subseteq A\}\) be the conductor, \(C=A/I\), \(D=B/I\). The conductor square gives \(A=B\times_D C\). With \(P_T=P\otimes_A T\), the aid reports the exact compatibility description
\[
\operatorname{Um}(P^\vee)\simeq
\operatorname{Um}(P_B^\vee)
\mathbin{\times}_{\operatorname{Um}(P_D^\vee)}
\operatorname{Um}(P_C^\vee).
\]
Thus a free summand downstairs requires compatible unimodular functionals, not merely separate summands on the normalization and conductor subscheme.

For \(A=k[t^2,t^3]\), it uses \(B=k[t]\), conductor \(I=(t^2,t^3)\), \(C=k\), \(D=k[t]/(t^2)\). Gluing trivial line bundles by \(1+at\in D^\times\) gives the aid's line bundle \(L_a\), nontrivial for each \(a\ne0\). It reports
\[
\mathrm{CH}_0(\operatorname{Spec}A)=0,
\qquad c_1(L_a)\cap[\operatorname{Spec}A]=0,
\]
in Fulton's homological Chow theory, although \(L_a\) has no free line summand, since a rank-one module with such a summand would be free. This supplies the aid's odd-dimensional, dimension-one counterexample to a direct singular extension of the preceding top-Chern detection question using that homological class. It does not refute every operational, motivic or other possible singular Chern-class formulation.
\subsection{Short English statement}
Ask which earlier projective-module and \(\A^1\)-homotopy questions extend to singular varieties and what is known about them. The attached aid treats free-summand descent and reports a cuspidal line-bundle example showing that vanishing of a Fulton homological top Chern class need not detect a free summand on a singular curve.
\subsection{Sources}
\begin{itemize}
\item Local UnsolvedMath registry, permanent ID \(20000247\), exact public code \texttt{catalogue entry}; canonical location \texttt{aim-algebraic-geometry-notes.json}, notes[246].
\item AIM workshop ``Projective modules and \(A^1\)-homotopy theory,'' Question 10, attributed in the supplied reconstruction to Anand Sawant; source URL \url{https://aimath.org/pastworkshops/projectiveA1problems.pdf}. The supplied verification says the PDF was inspected in its own run; this drafting step did not consult it. No year is supplied; the literature field is empty. The OCR tail ``4 SARANG SANE'' has unresolved page/header or attribution status.
\item Attached research attempt 1, ``Singular free-summand descent and a cuspidal top-Chern counterexample,'' referring to Fulton homological Chow groups without full bibliographic citations.
\end{itemize}
\end{document}
